\section{Структура HTML-отчёта: описание секций}
\label{app:html_report}

HTML-отчёт генерируется командой:
\begin{lstlisting}[language=sh]
python -m tools.ptx_report \
    --profiling profiling.txt \
    --ptx fft_kernel.ptx \
    --output report_fft.html
\end{lstlisting}

Отчёт содержит пять секций:

\textbf{Секция 1 --- Панель «пилюль»} (рисунок~\ref{fig:html_pills}):
сводные метрики ядра~--- общее число инструкций, число исполненных
потоко-операций, среднее $\overline{\eta}_{\mathrm{branch}}$,
суммарные $C_{\mathrm{bank}}$, средний $\bar{\kappa}$.

\begin{figure}[h]
  \centering
  \fbox{\parbox{12cm}{\centering\vspace{1.5cm}
    \textit{[TODO: скриншот секции панели пилюль для fft\_custom.\\
    Команда: cuemu --collect-profiling ./fft\_custom, затем python -m tools.ptx\_report,\\
    открыть HTML в браузере, сделать скриншот.]}
  \vspace{1.5cm}}}
  \caption{Панель «пилюль» HTML-отчёта.}
  \label{fig:html_pills}
\end{figure}

\textbf{Секция 2 --- Таблица инструкций} (рисунок~\ref{fig:html_table}):
все уникальные PC с тепловой подсветкой по \texttt{exec\_count};
столбцы: PC, мнемоника, блок, метрики.

\begin{figure}[h]
  \centering
  \fbox{\parbox{12cm}{\centering\vspace{1.5cm}
    \textit{[TODO: скриншот таблицы инструкций с тепловой подсветкой.]}
  \vspace{1.5cm}}}
  \caption{Таблица инструкций с тепловой подсветкой.}
  \label{fig:html_table}
\end{figure}

\textbf{Секция 3 --- Горячие точки}:
топ-10 инструкций по дивергенции (минимальный $\eta_{\mathrm{branch}}$),
по конфликтам банков (максимальный $C_{\mathrm{bank}}$),
по слиянию (минимальный $\kappa$).

\textbf{Секция 4 --- Таблица регистров}:
для каждого префикса регистров~--- тип, объявлено слотов,
накопленный \texttt{write\_ops}.

\textbf{Секция 5 --- Аннотированный исходный код}:
строки исходного \texttt{.cu}-файла с вложенными метриками
инструкций, привязанных через \texttt{.loc}-директивы.

\begin{figure}[h]
  \centering
  \fbox{\parbox{12cm}{\centering\vspace{1.5cm}
    \textit{[TODO: скриншот секции аннотированного исходного кода\\
    (fft\_custom.cu, секция butterfly-loop с метриками).]}
  \vspace{1.5cm}}}
  \caption{Аннотированный исходный код в HTML-отчёте.}
  \label{fig:html_source}
\end{figure}
